Advanced questions require four or more operations, often combining multi-condition filters, derived columns, grouped aggregations, and top-$k$ selection. Several involve thresholds that are computed from the data rather than supplied as constants.

\subsubsect{A1. Among entire home/apt listings in Manhattan available more than 200 days per year and priced under \$300, what are the top 3 neighbourhoods by average number of reviews?}

\textbf{Why multiple steps are required.}
Three independent filter conditions must all be satisfied before neighbourhood-level aggregation and ranking can be performed.

\textbf{Expected sequence of operations.}
\begin{enumerate}
  \item \texttt{filter\_rows}: \texttt{room\_type} is \texttt{"Entire home/apt"}
  \item \texttt{filter\_rows}: \texttt{neighbourhood\_group} is
    \texttt{"Manhattan"}
  \item \texttt{filter\_rows}: \texttt{availability\_365} $> 200$
  \item \texttt{filter\_rows}: \texttt{price} $< 300$
  \item \texttt{group\_and\_aggregate}: group by \texttt{neighbourhood}, compute
    \texttt{mean(number\_of\_reviews)} and \texttt{count}
  \item \texttt{sort\_rows}: sort by \texttt{mean\_reviews} descending
  \item \texttt{limit\_rows}: keep top 3
\end{enumerate}

\textbf{Classification.}
The depth of the filter chain (four conditions) combined with grouped aggregation and top-$k$ selection defines this as an advanced question.

% ------------------------------------------------------------------------------

\subsubsect{A2. Which neighbourhood has the highest concentration of ``high-value'' listings (where price \texorpdfstring{$>$}{is greater than} average price and availability \texorpdfstring{$>$}{is greater than} average availability and number of reviews \texorpdfstring{$>$}{is greater than} average number of reviews), among neighbourhoods with at least 30 total listings?}

\textbf{Why multiple steps are required.}
Each of the three thresholds must be computed from the data. A binary indicator must then be derived, grouped, and filtered by neighbourhood size before a concentration rate can be ranked.

\textbf{Expected sequence of operations.}
\begin{enumerate}
  \item \texttt{group\_and\_aggregate}: compute global means of \texttt{price},
    \texttt{availability\_365}, and \texttt{number\_of\_reviews}
  \item \texttt{derive\_columns}: compute \texttt{is\_high\_value = 1} if all
    three thresholds are exceeded, else \texttt{0}
  \item \texttt{group\_and\_aggregate}: group by \texttt{neighbourhood}, compute
    \texttt{count} as \texttt{total} and \texttt{sum(is\_high\_value)} as
    \texttt{hv\_count}
  \item \texttt{filter\_rows}: retain rows where \texttt{total} $\geq 30$
  \item \texttt{derive\_columns}: compute
    \texttt{hv\_rate = hv\_count / total}
  \item \texttt{sort\_rows}: sort by \texttt{hv\_rate} descending
  \item \texttt{limit\_rows}: keep top 1
\end{enumerate}

\textbf{Classification.}
This question requires two derivation steps, two aggregation steps, and data-driven thresholds for all three filter conditions -- the most compositionally complex pattern in the set.

% ------------------------------------------------------------------------------

\subsubsect{A3. What is the average price gap between shared room and private room listings within each borough, and which borough has the largest gap?}

\textbf{Why multiple steps are required.}
The price gap is not stored and cannot be computed in a single group-aggregate step. Per-room-type averages must first be computed, then the difference between two specific groups must be derived.

\textbf{Expected sequence of operations.}
\begin{enumerate}
  \item \texttt{filter\_rows}: retain rows where \texttt{room\_type} is
    \texttt{"Shared room"} or \texttt{"Private room"}
  \item \texttt{group\_and\_aggregate}: group by \texttt{neighbourhood\_group}
    and \texttt{room\_type}, compute \texttt{mean(price)}
  \item \texttt{derive\_columns}: pivot room type means into
    \texttt{avg\_shared} and \texttt{avg\_private} per borough, then
    compute \texttt{price\_gap = avg\_private - avg\_shared}
  \item \texttt{sort\_rows}: sort by \texttt{price\_gap} descending
  \item \texttt{limit\_rows}: keep top 1
\end{enumerate}

\textbf{Classification.}
The cross-group derivation step -- computing a difference between two aggregated values in the same table -- requires a restructuring operation that is absent from basic and intermediate questions.

% ------------------------------------------------------------------------------

\subsubsect{A4. Among listings with at least 20 reviews and a minimum stay of 1--3 nights, which listing in which borough has the lowest total cost (price \texorpdfstring{$\times$}{times} minimum\_nights)? Disregard listings with a price of \$0.}

\textbf{Why multiple steps are required.}
Four filter conditions narrow the dataset, a column must be derived, and the result must be sorted and limited to a single row.

\textbf{Expected sequence of operations.}
\begin{enumerate}
  \item \texttt{filter\_rows}: retain rows where \texttt{number\_of\_reviews}
    $\geq 20$
  \item \texttt{filter\_rows}: retain rows where \texttt{minimum\_nights}
    $\geq 1$ and \texttt{minimum\_nights} $\leq 3$
  \item \texttt{filter\_rows}: retain rows where \texttt{price} $> 0$
  \item \texttt{derive\_columns}: compute
    \texttt{total\_cost = price * minimum\_nights}
  \item \texttt{sort\_rows}: sort by \texttt{total\_cost} ascending
  \item \texttt{limit\_rows}: keep top 1
  \item \texttt{select\_columns}: retain listing name, borough, and
    \texttt{total\_cost}
\end{enumerate}

\textbf{Classification.}
The combination of four filter conditions, a derivation, and an ascending top-1 selection across the full listing-level table (rather than an aggregated group) distinguishes this from simpler questions.

% ------------------------------------------------------------------------------

\subsubsect{A5. Which hosts operating in all 5 boroughs have the highest average listing price, among hosts with at least 20 total listings?}

\textbf{Why multiple steps are required.}
The borough coverage of each host must be computed and compared to a threshold of 5 before any price ranking can be produced. This requires counting distinct values within a group -- a step that must precede the final filter and sort.

\textbf{Expected sequence of operations.}
\begin{enumerate}
  \item \texttt{filter\_rows}: retain rows where
    \texttt{calculated\_host\_listings\_count} $\geq 20$
  \item \texttt{group\_and\_aggregate}: group by \texttt{host\_name}, compute
    \texttt{count\_distinct(neighbourhood\_group)} as
    \texttt{borough\_count} and \texttt{mean(price)} as
    \texttt{avg\_price}
  \item \texttt{filter\_rows}: retain rows where \texttt{borough\_count}
    equals 5
  \item \texttt{sort\_rows}: sort by \texttt{avg\_price} descending
\end{enumerate}

\textbf{Classification.}
The requirement to count distinct categorical values within a group and use that count as a post-aggregation filter threshold is a sophisticated multi-step pattern that does not arise in basic or intermediate questions.